% anexos/anexoD_validacion.tex
% Los valores de SAP2000 vienen de validacion/resultados_sap2000.csv y las
% capturas, de capturas/: ver docs/05-validacion-sap2000.md.
\chapter{Validación con software comercial}\label{cap:validacion}

A lo largo del libro verificamos la calculadora con el cálculo a mano
(ejemplo A), con propiedades que deben cumplirse (simetría, equilibrio) y con
el método de las secciones (ejemplo C). Este anexo agrega una verificación
independiente: comparar sus resultados con los de SAP2000, un programa
comercial de análisis estructural de uso habitual en la práctica, para los
ejemplos B y C.

\avisovalidacion

% ---------------------------------------------------------------------------
\section{Cómo modelar una armadura en un programa comercial}\label{sec:d_modelar}

Los programas comerciales analizan pórticos, cuyas barras resisten flexión.
Para que reproduzcan el modelo de armadura del libro hay que cuidar cinco
cosas:
\begin{enumerate}
  \item \textbf{Barras articuladas.} Liberar los momentos en ambos extremos de
    cada barra, o analizar el modelo como armadura plana, sin giros en los
    nodos. Así cada barra trabaja solo con fuerza axial.
  \item \textbf{Sin peso propio.} Los programas suelen incluir el peso de las
    barras en las cargas; el libro no lo considera, de modo que hay que
    anularlo.
  \item \textbf{Las mismas propiedades.} Un material con
    $E = \SI{200}{\giga\pascal}$ y secciones definidas por su área, sin
    importar su forma.
  \item \textbf{Los mismos apoyos y cargas.} Un apoyo articulado en el nodo 1,
    uno de rodillo con el desplazamiento vertical restringido en el nodo 2
    (ejemplo B) o 7 (ejemplo C), y las cargas aplicadas en los nodos.
  \item \textbf{Unidades coherentes.} Trabajar en kilonewtons y metros, como
    la calculadora, o convertir los resultados antes de compararlos.
\end{enumerate}

SAP2000 dibuja los modelos en el plano $xz$: su eje $z$ es el vertical, que
en el libro llamamos $y$. Sus convenciones de signos coinciden con las del
libro: los desplazamientos y las reacciones son positivos en el sentido de
los ejes, y la fuerza axial $P$ es positiva en tracción.

Los resultados de los dos programas deben coincidir en sus primeras cifras
significativas. Diferencias pequeñas en la última cifra se deben al redondeo;
una diferencia mayor indica un error en uno de los dos modelos, casi siempre
en un apoyo, en las unidades o en el peso propio.

% ---------------------------------------------------------------------------
\section{Ejemplo B}\label{sec:d_ejemplo_b}

El ejemplo B (\cref{fig:ejemplo_b}) es hiperestático, de modo que sus fuerzas
dependen de la rigidez de las barras: es la prueba más exigente para la
calculadora. La \cref{fig:sap_b_modelo} muestra su modelo en SAP2000.

\begin{figure}[ht]
  \centering
  \captura{sap_b_modelo}{modelo del ejemplo B en SAP2000, con los números de
    los nodos y de las barras, los apoyos y las cargas.}
  \caption{El ejemplo B modelado en SAP2000.}
  \label{fig:sap_b_modelo}
\end{figure}

La \cref{tab:d_b_desplazamientos} compara los desplazamientos, y la
\cref{tab:d_b_fuerzas}, las reacciones y las fuerzas axiales.

\begin{table}[ht]
  \centering
  \caption{Ejemplo B: desplazamientos de los nodos, en milímetros.}
  \label{tab:d_b_desplazamientos}
  \small
  \begin{tabular}{@{}ccccc@{}}
    \toprule
    Nodo & $u_x$ calculadora & $u_x$ SAP2000 & $u_y$ calculadora & $u_y$ SAP2000 \\
    \midrule
    1 & \valor{d_b_ux_1} & \valorcomercial{d_b_ux_1} & \valor{d_b_uy_1} & \valorcomercial{d_b_uy_1} \\
    2 & \valor{d_b_ux_2} & \valorcomercial{d_b_ux_2} & \valor{d_b_uy_2} & \valorcomercial{d_b_uy_2} \\
    3 & \valor{d_b_ux_3} & \valorcomercial{d_b_ux_3} & \valor{d_b_uy_3} & \valorcomercial{d_b_uy_3} \\
    4 & \valor{d_b_ux_4} & \valorcomercial{d_b_ux_4} & \valor{d_b_uy_4} & \valorcomercial{d_b_uy_4} \\
    \bottomrule
  \end{tabular}
\end{table}

\begin{table}[ht]
  \centering
  \caption{Ejemplo B: reacciones y fuerzas axiales, en kilonewtons.}
  \label{tab:d_b_fuerzas}
  \small
  \begin{tabular}{@{}lcc@{}}
    \toprule
    Resultado & Calculadora & SAP2000 \\
    \midrule
    Reacción $x$ del nodo 1 & \valor{d_b_r_1} & \valorcomercial{d_b_r_1} \\
    Reacción $y$ del nodo 1 & \valor{d_b_r_2} & \valorcomercial{d_b_r_2} \\
    Reacción $y$ del nodo 2 & \valor{d_b_r_4} & \valorcomercial{d_b_r_4} \\
    \midrule
    Barra 1 & \valor{d_b_n_1} & \valorcomercial{d_b_n_1} \\
    Barra 2 & \valor{d_b_n_2} & \valorcomercial{d_b_n_2} \\
    Barra 3 & \valor{d_b_n_3} & \valorcomercial{d_b_n_3} \\
    Barra 4 & \valor{d_b_n_4} & \valorcomercial{d_b_n_4} \\
    Barra 5 & \valor{d_b_n_5} & \valorcomercial{d_b_n_5} \\
    Barra 6 & \valor{d_b_n_6} & \valorcomercial{d_b_n_6} \\
    \bottomrule
  \end{tabular}
\end{table}

Los resultados coinciden en todas las cifras que muestran las tablas. La
deformada de SAP2000 (\cref{fig:sap_b_deformada}) y su diagrama de fuerzas
axiales (\cref{fig:sap_b_axiales}) son, además, los mismos que dibuja nuestra
calculadora (\cref{fig:resultados_ejemplo_b}).

\begin{figure}[ht]
  \centering
  \captura{sap_b_deformada}{deformada del ejemplo B en SAP2000.}
  \caption{Deformada del ejemplo B en SAP2000.}
  \label{fig:sap_b_deformada}
\end{figure}

\begin{figure}[ht]
  \centering
  \captura{sap_b_axiales}{diagrama de fuerzas axiales del ejemplo B en
    SAP2000, con el valor de cada barra.}
  \caption{Fuerzas axiales del ejemplo B en SAP2000.}
  \label{fig:sap_b_axiales}
\end{figure}

% ---------------------------------------------------------------------------
\section{Ejemplo C}\label{sec:d_ejemplo_c}

El ejemplo C (\cref{fig:ejemplo_c}) es isostático: sus fuerzas dependen solo
del equilibrio, y ya las verificamos con el método de las secciones en el
\cref{cap:mi_calculadora}. Compararlas con SAP2000 confirma que ambos modelos
tienen la misma geometría y las mismas cargas; la flecha, en cambio, depende
de las áreas y del módulo de elasticidad, y verifica las rigideces.

\begin{figure}[ht]
  \centering
  \captura{sap_c_modelo}{modelo del ejemplo C en SAP2000, con los números de
    los nodos y de las barras, los apoyos y las cargas.}
  \caption{El ejemplo C modelado en SAP2000.}
  \label{fig:sap_c_modelo}
\end{figure}

\begin{table}[ht]
  \centering
  \caption{Ejemplo C: flecha y reacciones.}
  \label{tab:d_c_flecha}
  \small
  \begin{tabular}{@{}lcc@{}}
    \toprule
    Resultado & Calculadora & SAP2000 \\
    \midrule
    Desplazamiento $u_y$ del nodo 4 (mm) & \valor{d_c_uy_4} & \valorcomercial{d_c_uy_4} \\
    Reacción $y$ del nodo 1 (kN) & \valor{d_c_r_2} & \valorcomercial{d_c_r_2} \\
    Reacción $y$ del nodo 7 (kN) & \valor{d_c_r_14} & \valorcomercial{d_c_r_14} \\
    \bottomrule
  \end{tabular}
\end{table}

\begin{table}[ht]
  \centering
  \caption{Ejemplo C: fuerzas axiales, en kilonewtons.}
  \label{tab:d_c_fuerzas}
  \small
  \begin{tabular}{@{}cccccc@{}}
    \toprule
    Barra & Calculadora & SAP2000 & Barra & Calculadora & SAP2000 \\
    \midrule
    1 & \valor{d_c_n_1} & \valorcomercial{d_c_n_1} & 12 & \valor{d_c_n_12} & \valorcomercial{d_c_n_12} \\
    2 & \valor{d_c_n_2} & \valorcomercial{d_c_n_2} & 13 & \valor{d_c_n_13} & \valorcomercial{d_c_n_13} \\
    3 & \valor{d_c_n_3} & \valorcomercial{d_c_n_3} & 14 & \valor{d_c_n_14} & \valorcomercial{d_c_n_14} \\
    4 & \valor{d_c_n_4} & \valorcomercial{d_c_n_4} & 15 & \valor{d_c_n_15} & \valorcomercial{d_c_n_15} \\
    5 & \valor{d_c_n_5} & \valorcomercial{d_c_n_5} & 16 & \valor{d_c_n_16} & \valorcomercial{d_c_n_16} \\
    6 & \valor{d_c_n_6} & \valorcomercial{d_c_n_6} & 17 & \valor{d_c_n_17} & \valorcomercial{d_c_n_17} \\
    7 & \valor{d_c_n_7} & \valorcomercial{d_c_n_7} & 18 & \valor{d_c_n_18} & \valorcomercial{d_c_n_18} \\
    8 & \valor{d_c_n_8} & \valorcomercial{d_c_n_8} & 19 & \valor{d_c_n_19} & \valorcomercial{d_c_n_19} \\
    9 & \valor{d_c_n_9} & \valorcomercial{d_c_n_9} & 20 & \valor{d_c_n_20} & \valorcomercial{d_c_n_20} \\
    10 & \valor{d_c_n_10} & \valorcomercial{d_c_n_10} & 21 & \valor{d_c_n_21} & \valorcomercial{d_c_n_21} \\
    11 & \valor{d_c_n_11} & \valorcomercial{d_c_n_11} & & & \\
    \bottomrule
  \end{tabular}
\end{table}

También aquí los resultados coinciden, y el diagrama de fuerzas axiales de
SAP2000 (\cref{fig:sap_c_axiales}) muestra el mismo patrón de tracciones y
compresiones que el de nuestra calculadora en el \cref{cap:mi_calculadora}.

\begin{figure}[ht]
  \centering
  \captura{sap_c_axiales}{diagrama de fuerzas axiales del ejemplo C en
    SAP2000, con el valor de cada barra.}
  \caption{Fuerzas axiales del ejemplo C en SAP2000.}
  \label{fig:sap_c_axiales}
\end{figure}

% ---------------------------------------------------------------------------
\section{Conclusión}\label{sec:d_conclusion}

Para una armadura hiperestática y para una isostática de 21 barras, la
calculadora que construimos en este libro reproduce los desplazamientos, las
reacciones y las fuerzas axiales de un programa comercial. No es casualidad:
ambos resuelven el mismo sistema $\mK\,\vU = \vF$ con el mismo método. La
diferencia es que ahora sabemos qué hace cada paso.
